% 合成 fixture：基于证据的学术去 AI 模式簇。
\chapter{模式簇边界}\label{chap:pattern-clusters}

% 样例 A：H-ING、H-PROMO、H-ATTR、H-SCOPE、H-OUTLOOK 应命中。
\paragraph{样例 A。}
行业报告认为该框架具有突破性。表~\ref{tab:main} 中的准确率有所提高，彰显了该框架从表征
学习到临床部署的变革性影响。尽管仍有这些挑战，该方法将开辟令人期待的新方向。

% 样例 B：H-PRED 应命中。
\paragraph{样例 B。}
表~\ref{tab:main} 作为消融结果的展示。

% 样例 C：H-TERM 应命中。
\paragraph{样例 C。}
所提出的编码器压缩输入。随后，该特征提取器将结果送入表征模块，并与潜在处理器联合优化；
正文没有定义这些名称是否指向同一实体。

% 样例 D：有证据的分析尾句和评价。
\paragraph{样例 D。}
在噪声测试条件下，准确率提高了 4.2 个百分点；这一表述对应表~\ref{tab:main} 的已报告结果
和预注册分析~\cite{smith2024}。

% 样例 E：可解析的具名归因。
\paragraph{样例 E。}
Smith 等~\cite{smith2024} 报告该效应在既定噪声模型下适用于数据集 A，但未声称可以推广到
该设置之外。

% 样例 F：范围精确的真实三项枚举。
\paragraph{样例 F。}
评估测量准确率、延迟和峰值内存；表~\ref{tab:main} 为所有基线报告了这三项指标。

% 样例 G：精确技术谓词和显式术语定义。
\paragraph{样例 G。}
潜变量 $z$ 表示后验均值。本文将编码器定义为特征提取器，将解码器定义为重建模块；后文沿用
这些已声明含义。

% 样例 H：诚实局限与具体 future test。
\paragraph{样例 H。}
评估仅覆盖数据集 A，因此跨域效应可能无法推广。后续将使用相同的预注册指标和噪声模型测试
数据集 B。
